% pdflatex-tate-math-oneshot.tex -- 一回実行契約の実測
% 判定（TATE-PASS/TATE-FAIL）:
%   EM       : everymath が行内数式1式ごとにちょうど1回（$ と \( で2式）
%   ED       : everydisplay が表示数式1式ごとにちょうど1回（\[ と $$ で2式）
%   MATHMC   : \mathmc 本体の組版が「サイズ種類数」だけ（display/text の
%              同一サイズ分岐は box+copy 共有で1回; script/scriptscript は
%              実サイズが必要なため別組版。4分岐→3組版）
%   MBOX     : \mbox 本体の組版が1回（box は \mathchoice 分岐を持たない）
%   AMSTEXT  : \text は amstext 自身の4分岐実装のまま（再実装しない;
%              コンパイルと counter が amstext の firstchoice ゲート通り）
%   MATHLEVEL: \end{document} 時点で mathlevel が 0 に戻る（開始時処理が
%              式ごとに1回だけ実行され、増減が対になっている）
\documentclass[pdflatex,tate,paper=a5]{jlreq}
\usepackage{amsmath}
\pagestyle{empty}
\makeatletter
\newcount\emcount  \newcount\edcount
\newcount\mcbodycount \newcount\mbbodycount
\begin{document}
\everymath{\global\advance\emcount1 }
\everydisplay{\global\advance\edcount1 }
漢$x$漢\(y\)漢
\[
  E=mc^2
\]
$$a+b$$
\ifnum\emcount=2
  \typeout{TATE-PASS: EM (\the\emcount, expected 2)}%
\else
  \typeout{TATE-FAIL: EM (\the\emcount, expected 2)}%
\fi
\ifnum\edcount=2
  \typeout{TATE-PASS: ED (\the\edcount, expected 2)}%
\else
  \typeout{TATE-FAIL: ED (\the\edcount, expected 2)}%
\fi
% \mathmc 本体: 3種サイズの組版 = サイズ種類数（10pt分岐は box+copy 共有）
\everymath{}
\def\mcbody{\global\advance\mcbodycount1 }
漢$\mathmc{\mcbody 漢}$漢
\ifnum\mcbodycount=3
  \typeout{TATE-PASS: MATHMC (\the\mcbodycount, expected 3 = distinct sizes)}%
\else
  \typeout{TATE-FAIL: MATHMC (\the\mcbodycount, expected 3)}%
\fi
% \mbox 本体: 1回
\def\mbbody{\global\advance\mbbodycount1 }
漢$\mbox{\mbbody 漢}$漢
\ifnum\mbbodycount=1
  \typeout{TATE-PASS: MBOX (\the\mbbodycount, expected 1)}%
\else
  \typeout{TATE-FAIL: MBOX (\the\mbbodycount, expected 1)}%
\fi
% amstext \text: コンパイルゲート + firstchoice 互換（counter は
% amstext のパッチ済み \stepcounter が1回に制限する; 本体は4分岐）
漢$\text{\mbbody 和文}$漢
\ifnum\mbbodycount=5
  \typeout{TATE-PASS: AMSTEXT (body fires per amstext branch; \the\mbbodycount total incl.\ mbox)}%
\else
  \typeout{TATE-FAIL: AMSTEXT (\the\mbbodycount)}%
\fi
\ifnum\@math@level=0
  \typeout{TATE-PASS: MATHLEVEL (\the\@math@level)}%
\else
  \typeout{TATE-FAIL: MATHLEVEL (\the\@math@level, expected 0)}%
\fi
\end{document}
